\centerline{\bf \Huge Усреднения в графах}

\begin{flushright}
        \scriptsize{Эпизод 17. Четвёртый кандидат.}
\end{flushright}

\begin{tcolorbox}[
    enhanced,
    breakable,
    colback=gray!10,
    colframe=black!70,
    boxrule=0.6pt,
    arc=2mm
]
\textbf{Усреднение.}
Многие подсчёты на графах построены на идее «усреднения».
Если среднее значение величины равно $k$, то:
\begin{itemize}
    \item есть ситуация, когда величина не больше $k$;
    \item есть ситуация, когда величина не меньше $k$;
    \item либо во всех ситуациях величина равна $k$, либо есть и ситуации, когда величина больше $k$, и ситуации, когда величина меньше $k$.
\end{itemize}
\end{tcolorbox}

\textbf{Замечание.}
Обычно решение можно оформить двумя способами:
\begin{itemize}
    \item в прямую сторону через оценку среднего значения;
    \item от противного через двойную оценку суммы величин по всем ситуациям.
\end{itemize}

\begin{tcolorbox}[
    enhanced,
    breakable,
    colback=gray!10,
    colframe=black!70,
    boxrule=0.6pt,
    arc=2mm
]
\textbf{Важно.}
Усреднение даёт грубую оценку (как отношение площадей в клетчатых задачах), часто это лишь первый шаг к решению!
\end{tcolorbox}

\textbf{Свойство.}
$AM(x_1,\ldots,x_n)$ — среднее арифметическое чисел $x_1,\ldots,x_n$.
Если $a_1,\ldots,a_n$ и $b_1,\ldots,b_n$ — некоторые наборы чисел, $k$ и $l$ — константы, то
\begin{itemize}
    \item $AM(a_1+b_1,a_2+b_2,\ldots,a_n+b_n)=AM(a_1,a_2,\ldots,a_n)+AM(b_1,b_2,\ldots,b_n)$;
    
    \item $AM(ka_1+l,ka_2+l,\ldots,ka_n+l)=k\cdot AM(a_1,a_2,\ldots,a_n)+l$.
\end{itemize}

Распространённые средние в графах связаны с рёбрами:
\begin{itemize}
    \item средняя степень вершины;
    \item среднее число рёбер внутри некоторой конфигурации;
    \item среднее число общих знакомых (подсчёт «осьминожек»).
\end{itemize}

\begin{tcolorbox}[
    enhanced,
    breakable,
    colback=gray!10,
    colframe=black!70,
    boxrule=0.6pt,
    arc=2mm
]
\textbf{Пример.}
На шахматной доске отметили $25$ клеток. Докажите, что можно расставить $8$ не бьющих друг друга ладей так, чтобы не менее $4$ ладей стояли на отмеченных клетках.
\end{tcolorbox}

\begin{tcolorbox}[
    enhanced,
    breakable,
    colback=gray!10,
    colframe=black!70,
    boxrule=0.6pt,
    arc=2mm
]

\newpage

\hspace{3mm}

\textbf{Пример.}
\textbf{(а)} В компании из $20$ человек у каждого хотя бы $10$ знакомых.
Докажите, что найдутся такие две тройки людей, что любые два человека из разных троек знакомы.

\textbf{(б)} Какое минимальное количество знакомых должно быть у каждого, чтобы при таком способе решения гарантировать наличие $K_{2,3}$?
\end{tcolorbox}

\begin{enumerate}[label=\textbf{\arabic*.}, leftmargin=1.2cm]

    \item Телефонная компания ввела льготный тариф для школьников, позволяющий каждому школьнику выбрать $10$ человек, которым он может звонить бесплатно. Какое наибольшее количество школьников можно подключить к этому тарифу так, чтобы среди каждых двух школьников один мог бесплатно звонить другому?

    \item В классе $34$ ребёнка, некоторые из которых дружат (дружба взаимна). Каждый день в $2025$ году преподаватель собирал команду на матбой из $6$ детей, которые попарно дружат. Известно, что никакой состав не повторился более $3$ раз. Докажите, что есть ребёнок, который дружит хотя бы с $8$ одноклассниками.

    \item Докажите, что в любом графе можно удалить не более половины рёбер так, чтобы он стал двудольным.

    \item В деревне Мартышкино у каждого мальчика все знакомые с ним девочки знакомы между собой. У каждой девочки среди её знакомых мальчиков больше, чем девочек. Докажите, что в Мартышкино мальчиков живёт не меньше, чем девочек.

    \item В множестве $\{1,2,\ldots,100\}$ выбрано $45$ подмножеств так, что любой элемент лежит ровно в $36$ подмножествах. Докажите, что найдутся три подмножества, в объединении дающие всё множество.

    \item На улице стоят $103$ фонаря. В комнате управления находятся $322$ тумблера. Каждый фонарь соединён с $161$ тумблером. Никакие два тумблера не соединены с одним и тем же набором фонарей. Переключение тумблера меняет состояние всех фонарей, с которыми он соединён (горящие потухают и наоборот). Изначально все фонари зажжены. Докажите, что можно нажать на некоторые $67$ тумблеров таким образом, чтобы включилось не менее $52$ фонарей.

\end{enumerate}